\chapter{Metodología}

En este capítulo se describe la metodología seguida para el desarrollo del proyecto, detallando la organización del trabajo y el proceso de elaboración, revisión y retroalimentación que ha guiado la evolución tanto del documento como del videojuego desarrollado. La metodología adoptada ha condicionado de forma directa las decisiones de diseño, la planificación de tareas y el alcance final del proyecto.

El desarrollo del proyecto \textit{NeuroForge} se ha llevado a cabo siguiendo una metodología de carácter iterativo \cite{DesarrolloIterativo}, basada en la elaboración progresiva de versiones parciales que han sido evaluadas y pulidas de manera continua. Este enfoque se basa en la idea de que el proyecto evoluciona mediante ciclos sucesivos de trabajo, revisión y mejora, en los que cada iteración incorpora los aprendizajes y correcciones surgidos a partir de la anterior.

A diferencia de las metodologías clásicas de desarrollo de software orientadas a equipos numerosos y ciclos de producción cerrados, el presente proyecto se ha desarrollado en un contexto académico individual, lo que ha requerido un planteamiento metodológico más flexible. En este sentido, el trabajo se ha estructurado en una secuencia de entregas parciales en las que se han presentado avances tanto a nivel conceptual como técnico. Cada una de estas entregas ha sido revisada por los tutores del proyecto, quien ha proporcionado observaciones, correcciones y sugerencias que han servido como base para orientar el trabajo de la siguiente iteración.

Este proceso de revisión continua ha permitido una mejora progresiva del proyecto, ya que cada nueva versión del documento y del videojuego ha incorporado ajustes y ampliaciones respecto a la anterior. Además, la retroalimentación temprana ha facilitado la detección de errores conceptuales, problemas de diseño o decisiones técnicas poco adecuadas en fases iniciales del desarrollo, evitando así que dichos problemas se acumulen a etapas más avanzadas y complejas.

La metodología iterativa ha favorecido una adaptación constante al alcance real del Trabajo de Fin de Grado. A lo largo del desarrollo se han tenido en cuenta las limitaciones de tiempo, recursos y complejidad propias de un proyecto académico individual, permitiendo reajustar objetivos y prioridades de forma gradual. Este enfoque ha resultado especialmente relevante en el diseño del videojuego, donde las mecánicas, sistemas y elementos visuales han sido definidos, evaluados y refinados conforme se analizaba su viabilidad técnica y su coherencia con los objetivos del proyecto.

En la práctica, el desarrollo de \textit{NeuroForge} ha seguido una progresión estructurada que ha comenzado con la concepción de la idea y el diseño inicial del videojuego, pasando por la planificación del trabajo, el análisis de riesgos y la definición de requisitos, hasta culminar en la implementación de una versión completa y funcional. Esta evolución gradual ha permitido mantener un control continuo sobre el estado del proyecto y asegurar la coherencia entre los objetivos planteados inicialmente y los resultados obtenidos.

\section{Iteraciones del desarrollo}

Con el objetivo de organizar el trabajo de forma progresiva y controlada, el desarrollo del proyecto se estructuró en varias iteraciones. Cada iteración se centró en la implementación de un conjunto concreto de funcionalidades del videojuego, permitiendo validar el funcionamiento del sistema antes de avanzar hacia etapas de mayor complejidad o con funcionalidades dependientes de las de anteriores iteraciones.

\subsection{Iteración 1: Implementación del juego base}

La primera iteración tuvo como objetivo el desarrollo de una versión básica y completamente funcional del videojuego. En esta fase se implementaron las mecánicas fundamentales necesarias para que una partida pudiera desarrollarse correctamente sin necesidad de un bot, estableciendo las bases estructurales del sistema.

Durante esta iteración se definió en primer lugar la estructura del tablero y de las piezas que intervienen en el juego, lo que permitió establecer el marco sobre el que se desarrollaría el resto de la lógica del sistema. A continuación, se implementó el sistema de movimiento de las piezas, incluyendo las restricciones y reglas que determinan cómo pueden desplazarse dentro del tablero.

Una vez establecido el sistema de movimiento, se desarrolló el sistema de combate entre piezas, encargado de gestionar las interacciones que se producen cuando dos piezas entran en conflicto dentro del tablero. Paralelamente, se implementó el sistema de gestión del estado del tablero, que permite actualizar de forma correcta la posición de las piezas tras cada acción realizada por los jugadores.

Asimismo, se incorporó el sistema de turnos que regula la alternancia entre jugadores durante la partida. Finalmente, se configuró el estado inicial del juego al comenzar una nueva partida y se implementaron las condiciones que determinan la finalización de la misma. Durante todo este proceso se realizaron pruebas funcionales básicas con el objetivo de verificar el correcto funcionamiento de las mecánicas implementadas.

Al finalizar esta iteración se obtuvo una primera versión completamente jugable del juego, que incluía todas las mecánicas esenciales necesarias para el desarrollo de una partida.

\subsection{Iteración 2: Implementación de la inteligencia del bot}

Una vez establecida la base funcional del videojuego, la segunda iteración se centró en el desarrollo de un sistema de inteligencia artificial que permitiera controlar un oponente automático.

El objetivo principal de esta fase fue dotar al juego de un adversario capaz de analizar el estado del tablero y tomar decisiones de forma autónoma, respetando las reglas establecidas en la iteración anterior. Para ello, el desarrollo de la inteligencia del bot se planteó de forma progresiva, comenzando con un sistema de toma de decisiones básico y evolucionando posteriormente hacia un comportamiento más avanzado.

En una primera etapa se diseñó un bot capaz de identificar los movimientos válidos disponibles en cada turno. A partir de este conjunto de acciones posibles, el sistema selecciona una de ellas de forma parcialmente aleatoria. Este enfoque permite garantizar que el bot respeta las reglas del juego y es capaz de participar correctamente en una partida, aunque sus decisiones no sigan todavía una estrategia compleja.

Una vez implementado este comportamiento inicial, se plantea la evolución del sistema hacia un modelo de decisión más avanzado basado en el aprendizaje mediante repetición. En este enfoque, el bot puede entrenarse jugando múltiples partidas contra sí mismo, lo que permite ajustar progresivamente su comportamiento y mejorar la calidad de sus decisiones a lo largo del tiempo.

Durante esta iteración también se integró el sistema de inteligencia artificial dentro del sistema de turnos del juego, permitiendo que el bot pueda actuar automáticamente cuando corresponde su turno. Asimismo, se realizaron distintas pruebas de funcionamiento para analizar su comportamiento en diversas situaciones de partida y comprobar su correcta integración con el resto de mecánicas del sistema.

Tras completar esta iteración, el videojuego pasó a contar con una experiencia de juego funcional para un solo jugador, sentando además las bases para futuras mejoras en el comportamiento del bot.

\subsection{Iteración 3: Desarrollo de la interfaz de usuario}

La tercera iteración estuvo orientada a mejorar la interacción entre el jugador y el sistema mediante el desarrollo de una interfaz de usuario más clara e intuitiva.

Durante esta fase se comenzó con el proceso de diseño de la interfaz mediante la elaboración de distintos bocetos y propuestas visuales. Estos bocetos permitieron definir la disposición de los elementos principales de la pantalla, como el tablero, las piezas, los indicadores de turno y los diferentes elementos de interacción con el jugador. Este proceso de diseño previo facilitó la planificación de la interfaz antes de su implementación dentro del juego.

A partir de este diseño inicial se trabajó en la representación visual del tablero y en la representación gráfica de las piezas, con el objetivo de facilitar la comprensión del estado del juego. Sobre esta base visual se implementaron distintos indicadores que permiten al jugador identificar de forma clara qué piezas están seleccionadas y cuáles son los movimientos disponibles en cada momento.

Posteriormente, se desarrollaron distintos elementos de interacción que permiten gestionar la partida de manera más cómoda, como menús básicos para iniciar o reiniciar el juego. Paralelamente, se introdujeron mejoras en la visualización del estado de la partida, permitiendo mostrar información relevante como el turno actual o el resultado de determinadas acciones dentro del juego.

Con estas mejoras, esta iteración permitió aumentar significativamente la claridad visual del sistema y facilitar la interacción del usuario con el juego.

\subsection{Iteración 4: Mejoras audiovisuales y pulido final}

La última iteración se centró en la incorporación de mejoras visuales y sonoras, así como en la optimización general del sistema con el objetivo de obtener una versión final más completa y estable del videojuego.

Durante esta fase se añadieron efectos de sonido asociados a distintas acciones del juego, como movimientos o interacciones entre piezas, lo que contribuyó a mejorar la inmersión del jugador en la experiencia de juego. Además, se realizaron diversas mejoras visuales en el tablero y en la representación de las piezas, refinando su apariencia y coherencia estética.

Asimismo, se incorporaron animaciones asociadas a determinadas acciones del juego, como el desplazamiento de las piezas o las interacciones entre ellas. Estas animaciones permiten representar de forma más clara los cambios que se producen en el tablero y aportan mayor dinamismo a la experiencia de juego.

También se llevaron a cabo tareas de optimización del rendimiento del sistema con el fin de garantizar un funcionamiento fluido del videojuego. Finalmente, se realizaron pruebas adicionales destinadas a detectar posibles errores o comportamientos inesperados, lo que permitió corregir fallos y asegurar la estabilidad del sistema.

Esta última iteración permitió consolidar todos los elementos desarrollados en fases anteriores y obtener una versión final del videojuego preparada para su presentación en el contexto del TFG.